\section*{Bab \ref{ch:b3:x-ray-diffraction} --- Kristalografi dan Difraksi Sinar-X}

\begin{solution}{exo:b3:x-ray-diffraction:1}
$d = a/\sqrt N$: $d_{111} = \qty{325.66}{pm}$, $d_{200} = \qty{282.03}{pm}$, $d_{220} =
\qty{199.42}{pm}$.
\end{solution}

\begin{solution}{exo:b3:x-ray-diffraction:2}
FCC: pantulan-pantulan pertamanya adalah (111) dan (200).
$\sin\theta = \lambda\sqrt N/2a$: $2\theta =
\ang{43.32}$ dan \ang{50.45}.
\end{solution}

\begin{solution}{exo:b3:x-ray-diffraction:3}
P: ketujuhnya. I ($h + k + l$ genap): 110, 200, 211, 220. F (tidak campuran): 111, 200, 220.
\end{solution}

\begin{solution}{exo:b3:x-ray-diffraction:4}
$\mathbf a^* = \mathbf e_x/a$, $\mathbf b^* = \mathbf e_y/b$: kisi persegi panjang bersisi
$1/a$ dan $1/b$, dengan sisi panjang kisi langsung menjadi sisi pendek.
\end{solution}

\begin{solution}{exo:b3:x-ray-diffraction:5}
Rasio $\sin^2\theta$ $1 : 2 : 3 : 4 : 5$, yaitu $N = 2, 4, 6, 8, 10$: berpusat badan
(110, 200, 211, 220, 310). Dari garis pertama, $d_{110} = \lambda/2\sin(\ang{20.135}) =
\qty{223.77}{pm}$ dan $a = d\sqrt2 = \qty{316.5}{pm}$: tungsten.
\end{solution}

\begin{solution}{exo:b3:x-ray-diffraction:6}
$Z = \rho N_Aa^3/M = 1.99 \times 6.022\times10^{23} \times (6.2929\times10^{-8})^3/74.6 = 4.0$.
\end{solution}

\begin{solution}{exo:b3:x-ray-diffraction:7}
$\beta = \qty{0.40}{\degree} = \qty{6.98e-3}{rad}$, $\cos\theta = \cos\ang{15.85} = 0.962$:
$L = 0.9 \times 0.15406/(6.98\times10^{-3} \times 0.962) = \qty{21}{nm}$.
\end{solution}

\begin{solution}{exo:b3:x-ray-diffraction:8}
$F_{hkl} = f_{\ce{Cs+}} + f_{\ce{Cl-}}(-1)^{h+k+l}$: $F_{100} = 54 - 18 = 36$, $F_{110} = 72$:
intensitasnya berasio 1 : 4. Pusat sel ditempati ion yang berbeda dari ion di
sudut-sudutnya: kisinya primitif (kisi berpusat badan memerlukan isi yang
identik di kedua titik), dan 100 hadir.
\end{solution}

\begin{solution}{exo:b3:x-ray-diffraction:9}
Setiap atom mempunyai salinan yang tergeser sebesar $(\frac12,\frac12,0)$, yang
mengalikan $F$ dengan $1 +
\eu^{\iu\pi(h+k)} = 1 + (-1)^{h+k}$, nol untuk $h + k$ ganjil.
\end{solution}

\begin{solution}{exo:b3:x-ray-diffraction:10}
Kuasikristal mempunyai keteraturan jangka panjang tanpa periodisitas: kuasikristal
tidak mempunyai kisi, sehingga pembatasan itu, yang menyangkut kisi, tidak
berlaku baginya. Bintik-bintik difraksinya yang tajam membuat kristal
didefinisikan ulang sebagai setiap padatan yang pola difraksinya pada dasarnya
diskret, periodik atau tidak.
\end{solution}

\begin{solution}{exo:b3:x-ray-diffraction:11}
Luncuran $c$ memetakan $(x, y, z)$ ke $(x, \frac12 - y, \frac12 + z)$. Untuk $k = 0$
kedua faktor fasenya adalah $\eu^{2\pi\iu(hx + lz)}$ dan $\eu^{2\pi\iu(hx + lz + l/2)} = (-1)^l
\eu^{2\pi\iu(hx+lz)}$: jumlahnya lenyap untuk $l$ ganjil.
\end{solution}

\begin{solution}{exo:b3:x-ray-diffraction:12}
Inversi, cermin, atau luncuran mengubah sebuah molekul menjadi bayangan
cerminnya, sehingga kristal yang memuat salah satunya akan memuat kedua
enantiomer. Enantiomer murni hanya mengkristal dalam 65 \omterm{def:b3:x-ray-diffraction:space-group}{grup ruang} yang
tersusun dari rotasi, \omterm{def:b3:x-ray-diffraction:space-group}{sumbu ulir}, dan translasi (grup Sohncke), seperti
\textit{P}$2_12_12_1$. Menurut hukum Friedel, pola satu enantiomer sama
dengan pola enantiomer lainnya; konfigurasi absolut diperoleh dari
penyimpangan kecil yang ditimbulkan oleh hamburan anomali atom-atom yang
lebih berat.
\end{solution}

\begin{solution}{pb:b3:x-ray-diffraction:1}
\textbf{1.} $\theta = \ang{14.171}$, $d = \lambda/2\sin\theta = \qty{314.64}{pm}$.
\textbf{2.} 314.64, 222.49, 181.66, 157.32, 140.71, \qty{128.45}{pm}.
\textbf{3.} 1.000, 2.000, 3.000, 4.000, 5.000, 6.000.
\textbf{4.} 1, 2, 3, 4, 5, 6.
\textbf{5.} Ya, sebagai (100), (110), (111), (200), (210), (211) sel kubik
primitif dengan $a' = \qty{314.6}{pm}$.
\textbf{6.} $N = 7$ (dan 15, 23\dots), yang bukan jumlah tiga bilangan kuadrat;
perbedaan pertama baru akan muncul pada garis kedelapan. Enam garis tidak
dapat membedakan sel P bersisi $a'$ dari sel F bersisi $2a'$.
\textbf{7.} $N = 4, 8, 12, 16, 20, 24$: (200), (220), (222), (400), (420), (422).
\textbf{8.} $a = 2d_{200} = \qty{629.3}{pm}$.
\textbf{9.} \ce{KCl} (\qty{629.29}{pm}).
\textbf{10.} $F = 4[f(\ce{K+}) - f(\ce{Cl-})]$ untuk indeks yang semuanya ganjil, dan
kedua ion itu masing-masing mempunyai 18 elektron: amplitudonya saling
meniadakan.
\textbf{11.} Ya untuk keduanya: \ce{Na+} (10) dan \ce{Cl-} (18), \ce{K+} (18) dan
\ce{Br-} (36) menghamburkan secara berbeda; (111) \ce{NaCl} akan berada pada
\ang{27.36}, dan untuk \ce{KBr} pada \ang{23.38}, keduanya di dalam rentang yang
direkam.
\textbf{12.} Sinar-X melihat \ce{K+} dan \ce{Cl-} hampir identik; ion-ion identik
dalam susunan garam batu membentuk larik kubik sederhana bersisi $a/2$, sebuah
sel semu.
\textbf{13.} Empat \ce{KCl}.
\textbf{14.} $\rho = 4 \times 74.6/(6.022\times10^{23} \times (6.293\times10^{-8})^3) =
\qty{1.99}{g/cm^3}$.
\textbf{15.} Enam dan enam (oktahedron dari ion yang lain).
\textbf{16.} $a/2 = \qty{314.6}{pm}$.
\textbf{17.} (440), $N = 32$, pada \ang{87.65} (pantulan (333)/(511) yang semua
indeksnya ganjil tidak hadir).
\textbf{18.} Intensitas berubah karena orientasi kristalit yang disukai,
serapan, menurunnya faktor hamburan, dan gerak termal; posisi hanya
bergantung pada kisi.
\textbf{19.} Dari $d = \lambda/2\sin\theta$, $\Delta d/d = -\cot\theta\,\Delta\theta$: galat
sudut yang sama memberikan galat relatif yang lebih kecil pada $\theta$ yang
besar.
\textbf{20.} $d_{420} = \lambda/2\sin\ang{33.190} = \qty{140.71}{pm}$, $a = d\sqrt{20} =
\qty{629.3}{pm}$.
\textbf{21.} $L = 0.9 \times 0.15406/(4.36\times10^{-3} \times \cos\ang{33.19}) = \qty{38}{nm}$.
\textbf{22.} Tidak: $\beta = \ang{0.0095}$, jauh di bawah lebar instrumental yang
lazim.
\textbf{23.} \textbf{$a = \qty{629.3}{pm}$: serbuk itu kalium klorida}, yang
pantulan-pantulan ganjilnya dibungkam oleh dua ion dengan banyak elektron
yang sama.
\end{solution}
